% ── sections/05_references.tex ─────────────────────────────────────────────────
% Numbered reference list (Nature-style)

\begin{thebibliography}{99}

\bibitem{emgnn2023}
Chatzianastasis, M. \textit{et al.}
Explainable multilayer graph neural network for cancer gene prediction.
\textit{Bioinformatics} \textbf{39}, btad643 (2023).

\bibitem{emogi2021}
Schulte-Sasse, R. \textit{et al.}
Integration of multi-omics data with graph convolutional networks to identify
new cancer genes and their associated molecular mechanisms.
\textit{Nat. Mach. Intell.} \textbf{3}, 310--321 (2021).

\bibitem{gcn2017}
Kipf, T.\ N.\ \& Welling, M.
Semi-supervised classification with graph convolutional networks.
\textit{Proc. ICLR} (2017).

\bibitem{gat2018}
Veli\v{c}kovi\'{c}, P. \textit{et al.}
Graph attention networks.
\textit{Proc. ICLR} (2018).

\bibitem{gin2019}
Xu, K. \textit{et al.}
How powerful are graph neural networks?
\textit{Proc. ICLR} (2019).

\bibitem{sage2017}
Hamilton, W.\ L. \textit{et al.}
Inductive representation learning on large graphs.
\textit{Proc. NeurIPS} (2017).

\bibitem{ig2017}
Sundararajan, M. \textit{et al.}
Axiomatic attribution for deep networks.
\textit{Proc. ICML} (2017).

\bibitem{gsea2005}
Subramanian, A. \textit{et al.}
Gene set enrichment analysis: a knowledge-based approach for interpreting
genome-wide expression profiles.
\textit{Proc. Natl Acad. Sci. USA} \textbf{102}, 15545--15550 (2005).

\bibitem{enrichr2013}
Chen, E.\ Y. \textit{et al.}
Enrichr: interactive and collaborative HTML5 gene list enrichment analysis tool.
\textit{BMC Bioinformatics} \textbf{14}, 128 (2013).

\bibitem{msigdb2015}
Liberzon, A. \textit{et al.}
The Molecular Signatures Database (MSigDB) hallmark gene set collection.
\textit{Cell Syst.} \textbf{1}, 417--425 (2015).

\bibitem{hallmarks2011}
Hanahan, D.\ \& Weinberg, R.\ A.
Hallmarks of cancer: the next generation.
\textit{Cell} \textbf{144}, 646--674 (2011).

\bibitem{kingma2014adam}
Kingma, D.\ P.\ \& Ba, J.
Adam: a method for stochastic optimization.
\textit{Proc. ICLR} (2015).

\bibitem{batchnorm2015}
Ioffe, S.\ \& Szegedy, C.
Batch normalization: accelerating deep network training by reducing internal
covariate shift.
\textit{Proc. ICML} (2015).

\bibitem{resnet2016}
He, K. \textit{et al.}
Deep residual learning for image recognition.
\textit{Proc. CVPR} (2016).

\bibitem{cosine2017}
Loshchilov, I.\ \& Hutter, F.
SGDR: stochastic gradient descent with warm restarts.
\textit{Proc. ICLR} (2017).

\bibitem{graphnorm2021}
Cai, T. \textit{et al.}
GraphNorm: a principled approach to accelerating graph neural network training.
\textit{Proc. ICML} (2021).

\bibitem{cgc2022}
Wang, Y.\ \& Chatzianastasis, M.
Cancer gene classification using graph neural networks.
\textit{CS6230 Course Project, NUS} (2022).

\bibitem{captum2020}
Kokhlikyan, N. \textit{et al.}
Captum: a unified and generic model interpretability library for PyTorch.
\textit{arXiv preprint} arXiv:2009.07896 (2020).

\bibitem{lawrence2013}
Lawrence, M.\ S. \textit{et al.}
Mutational heterogeneity in cancer and the search for new cancer-associated genes.
\textit{Nature} \textbf{499}, 214--218 (2013).

\bibitem{focal2017}
Lin, T.-Y. \textit{et al.}
Focal loss for dense object detection.
\textit{Proc. ICCV} (2017).

\bibitem{dropedge2020}
Rong, Y. \textit{et al.}
DropEdge: towards deep graph convolutional networks on node classification.
\textit{Proc. ICLR} (2020).

\bibitem{graphmae2022}
Hou, Z. \textit{et al.}
GraphMAE: self-supervised masked graph autoencoders.
\textit{Proc. KDD} (2022).

\bibitem{gps2022}
Ramp\'{a}\v{s}ek, L. \textit{et al.}
Recipe for a general, powerful, scalable graph transformer.
\textit{Proc. NeurIPS} (2022).

\bibitem{hgnn2019}
Feng, Y. \textit{et al.}
Hypergraph neural networks.
\textit{Proc. AAAI} (2019).

\bibitem{pinnacle2024}
Li, M.\ M. \textit{et al.}
Contextual AI models for single-cell protein biology.
\textit{Nat. Methods} \textbf{21}, 1546--1557 (2024).

\bibitem{gnnexplainer2019}
Ying, R. \textit{et al.}
GNNExplainer: generating explanations for graph neural networks.
\textit{Proc. NeurIPS} (2019).

\bibitem{rwpe2020}
Dwivedi, V.\ P. \textit{et al.}
Benchmarking graph neural networks.
\textit{arXiv preprint} arXiv:2003.00982 (2020).

\bibitem{hipgnn2025}
Zang, Y. \textit{et al.}
Rethinking cancer gene identification through graph anomaly analysis.
\textit{Proc. AAAI} \textbf{39}, 13161--13169 (2025).

\bibitem{sgcd2024}
Li, X. \textit{et al.}
Towards simplified graph neural networks for identifying cancer driver genes in
heterophilic networks.
\textit{Brief. Bioinform.} \textbf{26}, bbae691 (2025).

\bibitem{dishyper2024}
Deng, C. \textit{et al.}
Identifying new cancer genes based on the integration of annotated gene sets and
gene networks using hypergraph neural networks.
\textit{Bioinformatics} \textbf{40}, i511--i520 (2024).

\bibitem{dghnn2025}
Li, Y. \textit{et al.}
DGHNN: a deep graph and hypergraph neural network for pan-cancer related gene
prediction.
\textit{Bioinformatics} \textbf{41}, btaf379 (2025).

\bibitem{deepcdg2025}
Wu, Y. \textit{et al.}
Deep graph convolutional network-based multi-omics integration for cancer driver
gene identification.
\textit{Brief. Bioinform.} \textbf{26}, bbaf364 (2025).

\bibitem{hanahan2022}
Hanahan, D.
Hallmarks of cancer: new dimensions.
\textit{Cancer Discov.} \textbf{12}, 31--46 (2022).

\end{thebibliography}
